\section{Formal Core}
\label{sec:theory}

Two conventions govern this section. Every guarantee is about what a finite
oracle certifies---$\Gpool$ where we can get it, $\Gk$ where the depth
matters---and never silently about $\Gstar$; and the
guard is analyzed as an \emph{index}, which changes the order in which stored
counterexamples are consulted but not the worst case, so we claim a measured
cost reduction and not an asymptotic separation. Statements and proofs of
everything below are in \appref{app:proofs}{the supplementary material}.

\begin{assumption}[Refutation soundness]
\label{as:sound}
When $\Ok$ refutes a candidate it returns a genuine witness: if
$\Ok(p) = x$ then $\lnot\spec(p, x)$.
\end{assumption}

\Cref{as:sound} constrains \emph{refutation}, not acceptance. It is the only
oracle property the guard needs; it holds for any oracle reporting a failing
execution it actually ran, and for ours it also requires the reference to be
right, which the corpus gate enforces (\cref{sec:design}). It has one soft
edge: a \emph{timeout} verdict is a resource fact rather than a witness of
$\lnot\spec$, and the one fault whose candidates sit at the sandbox limit
produces every cross-arm inconsistency we observe
(\cref{sec:results-integrity}). The converse---that acceptance implies
correctness---is assumed nowhere: at depth $\oracledepth$ it is false in
general, and \cref{sec:results-depth} measures how false.

\begin{theorem}[Guard soundness]
\label{thm:sound}
Under \cref{as:sound}, if $\mathrm{Guard}_\mem(p) = 1$ with witness $x$ then
$p \notin \Gpool$, and hence $p \notin \Gstar$; and $p \notin \Gk$ for every
check set containing $x$, which includes every $\Xk$ that exhausts the pool. No
candidate such a call would have accepted is ever blocked, however widely the
guard searches. If \CEGMem\ returns $\hat p$, then $\hat p$ passed every input
its accepting call drew.
\end{theorem}

The theorem is about the guard, not the oracle: \cref{eq:guard} blocks only on
a counterexample that is \emph{re-executed and still fails}, so widening the
search can only find more true refutations, and nothing here says
$\hat p \in \Gstar$. Its second clause can fail only on a redrawn $\Xk$ that
omits the witness, so we test rather than assume: an audit condition pays the
oracle on blocked rounds and asks whether it accepts, and over \num{4411}
blocked rounds it never did. That is not a tight bound on the hazard---where
$\Xk$ exhausts the pool the audit confirms what already holds by construction,
and below $\oracledepth = 100$, where the hazard is real, it was not run---but
it does establish that the implemented guard blocks nothing the implemented
oracle accepts. This is falsifier F1, with the requirement that a guarded arm's
repair rate equal its unguarded twin's.

\begin{corollary}[Budgeted success]
\label{cor:budget}
For any guard sound in the sense of \cref{thm:sound}: if blocked rounds are
\emph{not} charged against the proposal budget $\budget$, the probability of
repairing within $\budget$ charged proposals cannot fall, and rises strictly
only on realizations where a block moves the first accepting draw from beyond
$\budget$ to within it. If blocked rounds \emph{are} charged, the guard
is outcome-neutral: under common random numbers the round of first acceptance
is invariant to it.
\end{corollary}

\Cref{cor:budget} shaped the study and applies to untyped memory as much as
typed. Under the accounting \cref{sec:results} uses---one proposal, one budget
unit---the guard cannot change which round first accepts, so any gap between a
guarded arm and its unguarded twin is a soundness bug, not a result
(\cref{sec:results-integrity}); it also makes the first half untestable there,
hence the free-guarded condition (\cref{sec:results-ablation}).

\paragraph{Three further results, stated in full in the supplement.}
Under an i.i.d.\ proposer and \emph{ideal typing}---which also requires
refutation to be class-exclusive, an idealization our $\loc$ index does not
meet---typed memory halts on an unbounded run within $\Kreach{+}1$ oracle calls
\emph{and} $\Kreach{+}1$ proposals, $\Kreach$ being the reachable failure
classes, while untyped memory attains the call bound and no proposal bound.
Expected oracle calls fall from $1/\pi$ to $1 + \Dred$ with either memory,
$\Dred = \sum_\tau q_\tau/(q_\tau + \pi)$. And counting redundancy against the
\emph{observed-type history}---the classes already refuted this episode whether
or not the arm records them, which is what makes it measurable in a memory-free
arm---redundant attempts \emph{paid} fall to $0$ with either memory, while
those \emph{generated} fall to $0$ only with typed memory, which needs steering
to work; \cref{sec:results-ablation} reports that it does not. Of the three
budget quantities only two are bounded: not \emph{charged} proposals, which
include blocked rounds. These are directional predictions, and
\cref{sec:results-assumptions} reports which of the closed forms the data
supports.

\paragraph{What the index costs.}
Our guard consults the $\loc$-matched bucket first and then the rest of memory
(\cref{sec:impl}). Both it and a flat scan short-circuit at the first entry
that still refutes, and both consult all of memory when none does, so their
worst cases coincide and no asymptotic separation exists. What differs is where
in the order a refuting entry is found: the index saves to the extent that
refuting entries concentrate in the bucket---$35\%$ of guard decisions in vivo,
a rate the random-partition control matches (\cref{sec:results-typing})---and
pays $\lvert\mem[\loc(p)]\rvert$ extra consultations on a miss. On this corpus
there is little order to exploit: memory holds a median of one entry when the
guard fires, and $86\%$ of blocks resolve on the first consultation. Separately, deduplicating stored
counterexamples by input and memoizing verdicts within a round reduce
\emph{executions} at unchanged \emph{consultations}---which is the signature
\cref{sec:results-ablation} finds, and the reason our claim there is a measured
constant factor.
